%% ============================================================
%%  Centralized vs. Federated Learning — Comparison Report
%%  N-BaIoT IoT Anomaly Detection / 10–50 Client Scenarios
%% ============================================================
\documentclass[11pt,a4paper]{article}

\usepackage[margin=2.5cm]{geometry}
\usepackage{booktabs}
\usepackage{multirow}
\usepackage{array}
\usepackage{xcolor}
\usepackage{amsmath,amssymb}
\usepackage{hyperref}
\usepackage{caption}
\usepackage{subcaption}
\usepackage{enumitem}
\usepackage[expansion=false]{microtype}
\usepackage{parskip}
\usepackage{dcolumn}

\definecolor{accent}{HTML}{2563EB}

\hypersetup{colorlinks=true,linkcolor=accent,urlcolor=accent,citecolor=accent,
  pdftitle={Centralized vs. Federated Learning on N-BaIoT}}

\newcolumntype{N}{D{.}{.}{2}}

\begin{document}

% ---------------------------------------------------------------
\begin{titlepage}
  \centering\vspace*{3cm}
  {\LARGE\bfseries\color{accent}
    Centralized Training vs.\ Federated Learning\\[0.4em]
    for IoT Anomaly Detection\par}
  \vspace{1cm}
  {\large Comparison Report — N-BaIoT Dataset, 10–50 Client Scenarios\par}
  \vspace{0.5cm}
  {\normalsize\color{gray}
    Based on the federated learning framework proposed by\\
    Cao, Nicolau \& McDermott (2018)%
    \footnote{Cao, X., Nicolau, M., \& McDermott, J. (2018).
      \textit{Learning Neural Networks for Anomaly Detection
      in Cyber-Physical IoT Systems via Federated Learning}. Preprint.}
    and a centralized baseline built for direct comparison.\par}
  \vspace{2cm}{\large July 2026\par}\vfill
\end{titlepage}

% ---------------------------------------------------------------
\section*{Executive Summary}

This report compares \textbf{centralized learning} and \textbf{federated learning (FL)}
on the N-BaIoT benchmark across two experiments: per-client accuracy on 10 clients
(Experiment~1) and network-scale sensitivity from 10 to 50 clients (Experiment~3).
Both paradigms share identical model architectures, hyperparameters, data splits,
and evaluation metric (ROC-AUC).

Key findings:

\begin{itemize}[leftmargin=1.8em]
  \item \textbf{Centralized Autoencoder dominates on non-IID data at every scale},
        staying above $99.4\%$ AUC from 10 to 50 clients — a $+3$ to $+7$\,pp
        advantage over FL Autoencoders.
  \item \textbf{Centralized Autoencoder collapses on IID at 50 clients} ($95.23\%$),
        falling below every FL variant. FL Autoencoders degrade more gracefully
        ($97.0$–$99.1\%$ across scales).
  \item \textbf{FL Hybrid (SAE-CEN) is the most scale-stable method overall},
        holding $96.7$–$99.0\%$ on IID and $96.9$–$98.5\%$ on non-IID regardless
        of network size.
  \item \textbf{Centralized Hybrid underperforms FL Hybrid at every scale on
        non-IID data}, confirming that the centroid classifier benefits from
        per-client local statistics preserved during federation.
\end{itemize}

\newpage
% ---------------------------------------------------------------
\section{Background and Setup}

\subsection{Dataset}
N-BaIoT~\cite{meidan2018nbaiot} captures 115-dimensional network-flow statistics
from IoT gateways compromised by Mirai and BASHLITE botnets.
Two distribution regimes:
\textbf{IID} — clients share a homogeneous traffic profile;
\textbf{non-IID} — each client holds device-specific, heterogeneous traffic.

\subsection{Models}

\begin{description}[leftmargin=1.8em]
  \item[Autoencoder (AE)] MSE reconstruction loss; anomaly score = reconstruction error.
  \item[Hybrid (SAE-CEN)] Shrink Autoencoder with latent $\ell_2$ regularisation:
    $\mathcal{L} = \mathrm{MSE}(x,\hat{x}) + \frac{\lambda}{B}\sum_{j}\|z_j\|_2$,
    $\lambda=10$. Anomaly detection via centroid distance in latent space.
\end{description}

\subsection{Training Conditions}

\begin{center}
\small
\begin{tabular}{lllll}
\toprule
\textbf{Paradigm} & \textbf{Aggregation} & \textbf{Runs} & \textbf{Phases} & \textbf{Client ratio} \\
\midrule
Federated Learning & FedAvg, FedProx, MSE-Avg & 5 & 20 $\times$ 100 ep. & 50\% \\
Centralized        & —                         & 1 & 20 $\times$ 100 ep. & 100\% (merged) \\
\bottomrule
\end{tabular}
\end{center}

Data split identical in both: 40\% train / 10\% validation / 40\% dev (skipped
centralized) / 10\% test. One global StandardScaler in centralized vs.\ per-client
scaler in FL. Early stopping: patience = 3 epochs.

FL statistics for all scales are sourced from Cao et al.\ (2018); reuse is
legitimate since the centralized baseline uses the same code classes
(\texttt{Model}, \texttt{DataLoader}, \texttt{Evaluator}).

\newpage
% ---------------------------------------------------------------
\section{Experiment 1 — Per-Client Accuracy (10 Clients)}

\subsection{Aggregated results}

\begin{table}[ht]
\centering
\caption{Mean AUC (\%) on 10 clients. FL: mean $\pm$ std over 5 runs.
  Centralized: std over 10 clients at best phase.}
\label{tab:exp1-agg}
\small
\begin{tabular}{llrrrrrr}
\toprule
& & \multicolumn{3}{c}{\textbf{IID}} & \multicolumn{3}{c}{\textbf{non-IID}} \\
\cmidrule(lr){3-5}\cmidrule(lr){6-8}
\textbf{Paradigm} & \textbf{Variant} & Mean & Std & & Mean & Std & \\
\midrule
\multirow{6}{*}{FL}
  & AE + FedAvg      & \textbf{99.07} & 0.24 && 94.74 & 2.69 & \\
  & AE + FedProx     & 98.95          & 0.29 && 93.98 & 2.90 & \\
  & AE + MSE-Avg     & 98.92          & 0.31 && 92.87 & 1.26 & \\
  & SAE + FedAvg     & 98.76          & 0.56 && 96.93 & 0.70 & \\
  & SAE + FedProx    & 98.53          & 0.65 && \textbf{97.28} & 0.84 & \\
  & SAE + MSE-Avg    & 99.01          & 0.48 && 97.30 & 0.49 & \\
\midrule
\multirow{2}{*}{Centralized}
  & Autoencoder      & \textbf{99.78} & 0.12 && \textbf{99.71} & 0.08 & \\
  & Hybrid (SAE-CEN) & 97.67          & 0.29 && 95.40          & 1.84 & \\
\bottomrule
\end{tabular}
\end{table}

\subsection{Per-client breakdown — IID}

\begin{table}[ht]
\centering
\caption{Per-client AUC (\%) — IID, 10 clients. FL: averages over 5 runs.}
\label{tab:per-client-iid}
\tiny\setlength{\tabcolsep}{4pt}
\begin{tabular}{lrrrrrrrrrrr}
\toprule
\textbf{Method} & C1 & C2 & C3 & C4 & C5 & C6 & C7 & C8 & C9 & C10 & \textbf{Avg} \\
\midrule
\multicolumn{12}{l}{\textit{Federated Learning}} \\
AE + FedAvg     & 98.77 & 97.81 & 98.43 & 99.42 & 99.75 & 99.23 & 98.30 & 100.00 & 99.03 & 99.97 & \textbf{99.07} \\
AE + FedProx    & 98.53 & 97.53 & 98.10 & 99.39 & 99.75 & 99.16 & 98.08 & 100.00 & 98.94 & 99.97 & 98.95 \\
AE + MSE-Avg    & 98.54 & 97.51 & 98.02 & 99.32 & 99.73 & 99.09 & 98.06 & 100.00 & 98.93 & 99.97 & 98.92 \\
SAE + FedAvg    & 97.93 & 97.73 & 97.80 & 99.22 & 99.84 & 98.92 & 97.82 & 100.00 & 98.38 & 99.95 & 98.76 \\
SAE + FedProx   & 97.67 & 97.76 & 97.74 & 98.52 & 99.81 & 98.22 & 97.59 & 100.00 & 98.06 & 99.95 & 98.53 \\
SAE + MSE-Avg   & 98.68 & 97.83 & 98.26 & 99.31 & 99.86 & 99.16 & 98.30 & 100.00 & 98.72 & 99.95 & 99.01 \\
\midrule
\multicolumn{12}{l}{\textit{Centralized}} \\
Autoencoder     & 99.81 & 99.68 & 99.88 & 100.00 & 99.56 & 99.88 & 99.72 & 99.85 & 99.77 & 99.68 & \textbf{99.78} \\
Hybrid          & 97.80 & 97.50 & 97.83 & 97.79  & 97.02 & 97.73 & 97.48 & 97.90 & 98.12 & 97.52 & 97.67 \\
\bottomrule
\end{tabular}
\end{table}

\subsection{Per-client breakdown — non-IID}

\begin{table}[ht]
\centering
\caption{Per-client AUC (\%) — non-IID, 10 clients. FL: averages over 5 runs.}
\label{tab:per-client-noniid}
\tiny\setlength{\tabcolsep}{4pt}
\begin{tabular}{lrrrrrrrrrrr}
\toprule
\textbf{Method} & C1 & C2 & C3 & C4 & C5 & C6 & C7 & C8 & C9 & C10 & \textbf{Avg} \\
\midrule
\multicolumn{12}{l}{\textit{Federated Learning}} \\
AE + FedAvg     & 91.00 & 99.58 & 88.68 & 89.42 & 97.96 & 92.70 & 94.62 & 95.84 & 98.05 & 99.52 & 94.74 \\
AE + FedProx    & 87.96 & 99.58 & 87.70 & 88.72 & 97.53 & 91.75 & 93.94 & 95.47 & 97.65 & 99.52 & 93.98 \\
AE + MSE-Avg    & 86.61 & 99.56 & 86.80 & 86.14 & 96.91 & 89.47 & 92.53 & 94.18 & 97.04 & 99.48 & 92.87 \\
SAE + FedAvg    & 96.65 & 99.62 & 93.55 & 94.40 & 98.63 & 93.84 & 96.18 & 97.51 & 99.14 & 99.77 & 96.93 \\
SAE + FedProx   & 97.44 & 99.61 & 93.82 & 94.41 & 98.76 & 95.38 & 96.63 & 97.87 & 99.02 & 99.81 & \textbf{97.28} \\
SAE + MSE-Avg   & 97.25 & 99.64 & 94.58 & 94.97 & 98.69 & 94.43 & 96.58 & 97.73 & 99.30 & 99.77 & 97.30 \\
\midrule
\multicolumn{12}{l}{\textit{Centralized}} \\
Autoencoder     & 99.70 & 99.77 & 99.77 & 99.79 & 99.77 & 99.68 & 99.61 & 99.64 & 99.82 & 99.59 & \textbf{99.71} \\
Hybrid          & 93.71 & 95.41 & 94.84 & 96.46 & 97.40 & 94.29 & 95.77 & 97.35 & 97.42 & 91.34 & 95.40 \\
\bottomrule
\end{tabular}
\end{table}

\newpage
% ---------------------------------------------------------------
\section{Experiment 3 — Network Scale Sensitivity (10–50 Clients)}

Each cell shows mean AUC (\%).  For FL, only the best-performing aggregation
strategy per model family is shown.\footnote{FL best per scale:
  IID AE → FedAvg; IID SAE → MSE-Avg (10c), FedProx (20c), MSE-Avg (30–40c),
  FedAvg (50c).
  nonIID AE → MSE-Avg (30c, 50c), FedAvg otherwise; nonIID SAE → MSE-Avg (10–30c,
  50c), FedAvg (40c).}
$\Delta$ = Centralized $-$ FL best (positive favours centralized).

\begin{table}[ht]
\centering
\caption{Mean AUC (\%) by network scale — \textbf{IID} setting.}
\label{tab:scale-iid}
\small
\begin{tabular}{llrrrrr}
\toprule
\textbf{Paradigm} & \textbf{Model}
  & \textbf{10c} & \textbf{20c} & \textbf{30c} & \textbf{40c} & \textbf{50c} \\
\midrule
\multirow{2}{*}{Centralized}
  & Autoencoder  & \textbf{99.78} & \textbf{99.75} & \textbf{99.84} & \textbf{99.22} & 95.23 \\
  & Hybrid       & 97.67          & 97.42          & 97.19          & 96.84          & 96.87 \\
\midrule
\multirow{2}{*}{FL (best)}
  & AE           & 99.07          & 98.56          & 97.45          & 97.16          & \textbf{97.20} \\
  & SAE          & 99.01          & 99.02          & 98.34          & 98.45          & 98.33 \\
\midrule
\multirow{2}{*}{$\Delta$}
  & AE           & $+0.71$ & $+1.19$ & $+2.39$ & $+2.06$ & $-1.97$ \\
  & Hybrid/SAE   & $-1.34$ & $-1.60$ & $-1.15$ & $-1.61$ & $-1.46$ \\
\bottomrule
\end{tabular}
\end{table}

\begin{table}[ht]
\centering
\caption{Mean AUC (\%) by network scale — \textbf{non-IID} setting.}
\label{tab:scale-noniid}
\small
\begin{tabular}{llrrrrr}
\toprule
\textbf{Paradigm} & \textbf{Model}
  & \textbf{10c} & \textbf{20c} & \textbf{30c} & \textbf{40c} & \textbf{50c} \\
\midrule
\multirow{2}{*}{Centralized}
  & Autoencoder  & \textbf{99.71} & \textbf{99.83} & \textbf{99.81} & \textbf{99.79} & \textbf{99.49} \\
  & Hybrid       & 95.40          & 97.13          & 96.00          & 96.01          & 97.36 \\
\midrule
\multirow{2}{*}{FL (best)}
  & AE           & 94.74          & 95.72          & 96.98          & 95.99          & 96.42 \\
  & SAE          & 97.30          & 97.29          & 97.73          & \textbf{97.81} & \textbf{98.52} \\
\midrule
\multirow{2}{*}{$\Delta$}
  & AE           & $+4.97$ & $+4.11$ & $+2.83$ & $+3.80$ & $+3.07$ \\
  & Hybrid/SAE   & $-1.90$ & $-0.16$ & $-1.73$ & $-1.80$ & $-1.16$ \\
\bottomrule
\end{tabular}
\end{table}

\newpage
% ---------------------------------------------------------------
\section{Discussion}

\subsection{IID — Centralized Autoencoder breaks at scale}

On IID data, the centralized Autoencoder is the strongest method up to 40 clients
($99.22\%$), outperforming all FL variants by $+1$ to $+2.4$\,pp.
At 50 clients, however, it collapses to $95.23\%$ — falling $1.97$\,pp
\emph{below} the best FL variant (FedAvg at $97.01\%$).

This inversion is likely caused by \textbf{optimisation saturation}: the merged
training set at 50 clients is very large (around $33{,}000$ normal samples per
epoch), and the global model may overfit the dominant statistical modes in the
merged distribution while losing specificity for less-represented client profiles.
FL avoids this by keeping gradients local and aggregating only model parameters.

\subsection{non-IID — Centralized Autoencoder is scale-invariant}

On non-IID data the centralized Autoencoder stays remarkably stable: $99.71\%$ at
10 clients, $99.83\%$ at 20, $99.81\%$ at 30, $99.79\%$ at 40, and $99.49\%$ at
50 — a total drop of only $0.34$\,pp.  By contrast, FL Autoencoders hover between
$92.9\%$ and $97.0\%$ across scales.

The advantage is structural: by merging all device traffic, the centralized model
builds a comprehensive normal manifold that generalises across all device-specific
patterns.  Non-IID heterogeneity, which hurts FL aggregation, is a direct
\emph{benefit} in the centralized case.

\subsection{Hybrid (SAE-CEN) — FL wins at every scale on non-IID}

The centralized Hybrid underperforms its FL counterpart at every scale on non-IID
data (gap: $-0.16$ to $-1.90$\,pp).  The centroid classifier in the centralized
case is fitted on mixed latent codes from all devices, blurring device-specific
normal clusters.  In FL, the server broadcasts a shared encoder but each client
evaluates with its own training data for centroid placement, preserving local
anomaly thresholds.

On IID data, the FL Hybrid is also consistently better ($-1.15$ to $-1.61$\,pp).
The Hybrid's advantage from federation is thus structural, not distribution-dependent.

\subsection{Scale stability}

\begin{itemize}[leftmargin=1.8em]
  \item \textbf{Most stable overall}: FL Hybrid (SAE-CEN), regardless of
        distribution. Standard deviation across scales $\leq 0.5$\,pp.
  \item \textbf{Most stable centralized}: Centralized Autoencoder on non-IID.
        Essentially flat from 10 to 50 clients.
  \item \textbf{Least stable}: FL Autoencoder on non-IID (range: $92.87$ to $96.98\%$);
        Centralized Autoencoder on IID (range: $95.23$ to $99.84\%$).
\end{itemize}

\subsection{Deployment recommendation}

\begin{enumerate}[leftmargin=1.8em]
  \item \textbf{If data can be centralised and the network is small ($\leq$40 clients)}:
        use the Centralized Autoencoder — highest AUC, lowest variance, simplest
        infrastructure.
  \item \textbf{If data must remain local, or the network exceeds 40 clients}:
        use FL Hybrid (SAE-CEN + FedProx or MSE-Avg) — consistently above $96.7\%$
        on IID and $96.9\%$ on non-IID at all scales, with tight inter-run variance.
  \item \textbf{Avoid}: FL Autoencoder on non-IID at small scales ($<$30 clients)
        — AUC drops below $93\%$ and inter-run variance exceeds $2.9$\,pp.
\end{enumerate}

\newpage
% ---------------------------------------------------------------
\section{Conclusion}

\begin{enumerate}[leftmargin=1.8em]
  \item Centralized Autoencoder is the best-performing method on non-IID data at
        \emph{every} tested scale, outperforming FL by $+3$ to $+5$\,pp.
        It is, however, not viable beyond 40 clients on IID data.

  \item FL Hybrid (SAE-CEN) is the recommended strategy for large-scale or
        privacy-constrained deployments: its performance is insensitive to both
        distribution heterogeneity and network growth.

  \item The comparison validates using FL paper statistics directly alongside
        the centralized baseline, since both share the same model architecture,
        data pipeline, and evaluation protocol.

  \item A natural next step is experiment 2 (client participation ratio) which
        has no direct centralized mirror, but could be approximated by subsampling
        the training set to simulate partial data availability.
\end{enumerate}

% ---------------------------------------------------------------
\appendix
\section{Experimental Configuration}

\begin{center}
\small
\begin{tabular}{ll}
\toprule
\textbf{Parameter} & \textbf{Value} \\
\midrule
Dataset            & N-BaIoT (Mirai + BASHLITE) \\
Features           & 115 (network flow statistics) \\
Network scales     & 10, 20, 30, 40, 50 clients \\
Architectures      & Autoencoder, Shrink Autoencoder (SAE-CEN / Hybrid) \\
Latent dimension   & 11 \\
Hidden neurons     & 50 \\
Learning rate      & $10^{-5}$ \\
Batch size         & 12 \\
Shrinkage $\lambda$& 10 \\
Optimizer          & Adam \\
Early stopping     & patience = 3 epochs \\
Phases / rounds    & 20 \\
Epochs per phase   & 100 \\
FL runs            & 5 \\
FL client ratio    & 50\% \\
FL strategies      & FedAvg, FedProx ($\mu=0.01$), MSE-Avg \\
Data split         & 40\% train / 10\% val / 40\% dev / 10\% test \\
Scaler             & StandardScaler (per-client FL; global centralized) \\
Random seed        & 1234 \\
Metric             & ROC-AUC \\
\bottomrule
\end{tabular}
\end{center}

% ---------------------------------------------------------------
\begin{thebibliography}{9}

\bibitem{cao2018}
Cao, X., Nicolau, M., \& McDermott, J. (2018).
\textit{Learning Neural Networks for Anomaly Detection in Cyber-Physical IoT
Systems via Federated Learning}. Preprint.

\bibitem{meidan2018nbaiot}
Meidan, Y., et al. (2018).
N-BaIoT — Network-based Detection of IoT Botnet Attacks Using Deep Autoencoders.
\textit{IEEE Pervasive Computing}, 17(3), 12--22.

\bibitem{mcmahan2017}
McMahan, H. B., et al. (2017).
Communication-efficient Learning of Deep Networks from Decentralized Data.
\textit{AISTATS}.

\bibitem{li2020fedprox}
Li, T., et al. (2020).
Federated Optimization in Heterogeneous Networks. \textit{MLSys}.

\end{thebibliography}

\end{document}
